\documentclass[11pt]{article}
\usepackage[margin=1in]{geometry}
\usepackage{booktabs}
\usepackage{hyperref}
\usepackage{longtable}
\usepackage{amsmath}
\usepackage{xcolor}

\title{Replication Audit: Friedland, Kundr\'at \& Jacob (2012) --- Stochastic
Modelling of DSB Repair after Photon and Ion Irradiation}
\author{Ollie subagent (backfill 2026-07-05) --- LUCID-100 slot 67}
\date{Backfilled 2026-07-06 CDT}

\begin{document}
\maketitle

\begin{abstract}
This is a backfill audit of an existing LUCID-100 replication slot. The queue
labelled this row \textbf{REPLICATED}; the on-disk evidence supports
\textbf{NO-GO} (previously re-tiered from SPOT-CHECK on 2026-06-25).  We
preserve the honest on-disk verdict (\emph{NO-GO}) and flag the queue mismatch
in \texttt{failure\_analysis.md}. The paper is closed-access (Taylor \& Francis
IJRB, DOI 10.3109/09553002.2011.611404); its underlying Monte Carlo code
PARTRAC is proprietary (Helmholtz Zentrum M\"unchen, no public release); its
precursor parameter tables (Friedland 2010 RR1965) and calibration kinetics
(Stenerl\"ow 2000) are also closed. What exists in this slot is an independent
\emph{analytical reduction} of the paper's stated model refinements, not a
re-implementation of the PARTRAC stochastic Monte Carlo. This is a hard
reproducibility ceiling.
\end{abstract}

\section{Paper Summary (as far as we can reconstruct without the PDF)}

Friedland, Kundr\'at \& Jacob (2012), \emph{Int.\ J.\ Radiat.\ Biol.}
88(1--2):129--136, refines the NHEJ repair sub-model of the PARTRAC
track-structure Monte Carlo to better reproduce measured DSB rejoining
kinetics after low-LET ($^{60}$Co $\gamma$) and high-LET (nitrogen ion,
$\sim$80 keV/$\mu$m) irradiation. From the Semantic Scholar TLDR --- the only
paper text that could be quoted, since the S2 abstract field is elided by the
publisher --- the two headline refinements are (i) an \emph{ongoing
production of detectable DSB in the initial phase} (labile-site delayed
detection), and (ii) a slow-channel component with LET-dependent capacity
(complex-DSB rejoining is rate-limited). PARTRAC then simulates the
stochastic Gillespie-style NHEJ dynamics on DSB ends generated by the
microdosimetric track-structure module. The paper compares model output to
measured rejoining curves (Stenerl\"ow 2000 for photons + N-ions; other
sources for auxiliary species/energies).

\section{Claims Table}

\begin{longtable}{@{}p{0.4cm}p{7.5cm}p{2cm}p{2cm}p{2cm}@{}}
\toprule
\# & Claim (inferred from TLDR + reference graph) & Type & Testable here? & Tested? \\
\midrule
\endhead
C1 & Two-component (fast+slow) rejoining suffices to describe both photon and ion data with one model class & mechanistic & Directionally & Yes \\
C2 & $^{60}$Co $\gamma$ fast component dominates; fast $t_{1/2}\sim$~min; slow $t_{1/2}\sim$~h & quantitative range & Directionally & Yes \\
C3 & High-LET ions have larger slow (complex-DSB) fraction than photons & directional & Yes & Yes \\
C4 & High-LET ions have larger long-term (24~h) residual unrejoined fraction & directional & Yes & Yes \\
C5 & An ``ongoing production of detectable DSB in the initial phase'' (labile-site processing) improves fit vs a pure biexponential & mechanistic & Directionally & Yes \\
C6 & Same parameter set explains photon+ion with only an LET-dependent complex-DSB fraction & mechanistic & Weakly (via Hill saturation surrogate) & Yes \\
--- & Specific PARTRAC rate constants (fast/slow $k$, labile $k_{\text{lab}}$, capacity, etc.) & quantitative & \textbf{No} (closed) & No \\
--- & Specific DSB-complexity distributions per LET bin & quantitative & \textbf{No} (closed) & No \\
--- & Specific figure-level RMSE vs Stenerl\"ow 2000 data & quantitative & \textbf{No} (closed) & No \\
\bottomrule
\end{longtable}

\section{Method}
\begin{enumerate}
\item Resolved provenance via S2, OpenAlex, Unpaywall (2026-06-09). Confirmed closed access, no OA copy, no preprint, S2 abstract elided.
\item Mapped 14-reference graph (\texttt{source/references\_table.md}).
\item Searched GitHub/Zenodo/Google for public PARTRAC release: \textbf{none found} (only unrelated ``partrac''-named utilities).
\item Wrote \texttt{code/smoke\_friedland2012.py}: analytical model
  $F(t) = f\,e^{-k_f t} + (1{-}f)\,e^{-k_s t} + A_{\text{lab}}(1-e^{-k_{\text{lab}}t})\,e^{-k_s t}$,
  fitted independently to literature-typical $^{60}$Co $\gamma$ and N-ion reference curves by \texttt{scipy.optimize.curve\_fit} (5 params/regime, bounded NLS). Deterministic, no RNG, sub-second runtime.
\item Wrote \texttt{code/let\_sweep\_friedland2012.py}: LET sweep over 0.3, 1, 3, 10, 30, 60, 100, 150 keV/$\mu$m using a Hill saturation surrogate $f_{\text{complex}}(\text{LET}) = 0.55\,\text{LET}^{1.6}/(40^{1.6}+\text{LET}^{1.6})$ and $k_{\text{slow,eff}} = k_{\text{slow,base}}\,(1 - 0.85\,f_{\text{complex}})$.
\item Reran smoke to confirm bit-for-bit reproducibility.
\end{enumerate}

No paid endpoints; no HPC; CherryRd local Python 3 + numpy/scipy/matplotlib.

\section{Results vs Paper}

\begin{tabular}{@{}p{6.5cm}p{2cm}p{2cm}p{3.5cm}@{}}
\toprule
Endpoint & $^{60}$Co $\gamma$ & N-ion & Paper value \\
\midrule
Fast fraction $f$ & 0.76 & 0.63 & unknown (closed) \\
Fast $t_{1/2}$ (min) & 19.1 & 34.7 & unknown (closed) \\
Slow $t_{1/2}$ (min) & 398 & 1156 & unknown (closed) \\
Labile amp $A_{\text{lab}}$ & 0.017 & 0.029 & unknown (closed) \\
Residual @ 24~h (model) & 0.021 & 0.16 & unknown (closed) \\
Residual @ 24~h (input curve) & 0.06 & 0.18 & unknown (closed) \\
Fit RMSE & 0.021 & 0.025 & unknown (closed) \\
\bottomrule
\end{tabular}

\vspace{1em}

Smoke checks S1--S6: \textbf{6/6 pass}. LET-sweep trend checks T1--T6:
\textbf{4/6 pass} (T5 fails: high-LET residual $>$~15\% at 150 keV/$\mu$m ---
we get 12.2\%; T6 fails: high-LET slow $t_{1/2}$ $>$3$\times$ low-LET ---
Hill saturation capped at 0.85 keeps ratio $\sim$3$\times$, not $\ge$3$\times$).

\section{Per-Claim What Worked / What Didn't}
\begin{itemize}
\item \textbf{C1--C4 (directional):} recovered qualitatively by the analytical reduction.  This is a low bar: any two-component biexponential with photon-favoured fast component will hit these.
\item \textbf{C5 (labile-site):} $A_{\text{lab}}$ optimises to nonzero positive values for both regimes, consistent with the paper's own conclusion that the labile term improves fit --- but it \emph{cannot} be shown to be uniquely required without a formal AIC/BIC comparison versus the pure biexponential, which we did not perform.
\item \textbf{C6 (single parameter set + LET-dependent complexity):} very weakly supported. Our surrogate is a hand-tuned Hill sigmoid on $f_{\text{complex}}(\text{LET})$; it produces monotone trends but the T5/T6 magnitude checks fail. The paper's actual mechanism is PARTRAC track-structure-derived complexity distributions per LET bin, which we cannot reproduce.
\item \textbf{Quantitative claims:} \emph{not tested at all.} Zero paper-specific numbers were verified. This is the critical failure.
\end{itemize}

\section{Critique of the Replication (honest)}

This slot is far weaker than a queue label ``REPLICATED'' would suggest. A
genuinely honest self-assessment:
\begin{enumerate}
\item \textbf{PARTRAC was not re-implemented.} Not partially, not as a
Gillespie surrogate --- \emph{not at all}. What was written is an analytical
five-parameter biexponential-plus-labile-term, fitted to
\emph{literature-typical} rejoining curves, not the paper's actual data. The
paper's contribution is a stochastic track-structure Monte Carlo; ours is a
closed-form curve fit. These are different scientific objects.
\item \textbf{Stochastic distributions were not reproduced.} The paper's
signature output is a \emph{distribution} of DSB complexity and repair time
per track --- variance, tails, event-by-event fluctuations. We reproduced
only \emph{mean} curves. No second moment, no tail, no per-event
statistics. The word ``stochastic'' in the paper title does not survive
into our audit.
\item \textbf{Ion coverage is a single point.} The paper's story is
photon-vs-ion \emph{across a species/LET grid}. We tested one high-LET point
(N-ion $\sim$80 keV/$\mu$m). The LET sweep is a surrogate Hill sigmoid, not
a species-specific PARTRAC scan; it neither uses real ion physics nor is
calibrated to any species-specific measurement.
\item \textbf{Reference curves are unverified digitisations.} The
``literature-typical'' curves in \texttt{code/smoke\_friedland2012.py} lines
64--75 are hand-authored to look like Stenerl\"ow 2000 --- but nobody
verified them against the actual Stenerl\"ow paper (which is also closed).
The fit RMSE numbers (0.021, 0.025) are agreement-with-our-own-inputs, not
agreement-with-the-paper.
\item \textbf{Six/six smoke checks passing is close to tautological.} S1--S6
are chosen to accept the model class that was fit; passing them is
essentially a self-consistency check on the fit, not evidence for the
paper.
\item \textbf{Verdict was already reclassified NO-GO.} On 2026-06-25 the top
of REPORT.md was re-tagged from SPOT-CHECK to NO-GO under Rick's
hard-ceiling rule (nothing quantitatively reproducible $\Rightarrow$
NO-GO). The current backfill preserves that on-disk verdict.  The queue
label REPLICATED is stale/incorrect.
\end{enumerate}

\section{Verdict + Justification}

\textbf{VERDICT = NO-GO} (preserved from the 2026-06-25 re-tier).

\textbf{Queue mismatch:} the LUCID queue verdict for slot 67 is REPLICATED.
The on-disk evidence supports NO-GO. This mismatch is flagged in
\texttt{failure\_analysis.md}. \emph{Do not silently upgrade the on-disk
verdict to match the queue.}

Justification: PARTRAC is proprietary and unavailable; the paper PDF is
closed and its specific numbers are unreadable; the calibration data
(Stenerl\"ow 2000) and the precursor parameter tables (Friedland 2010
RR1965) are both closed. What we produced is a qualitative analytical
surrogate of the paper's \emph{stated} model refinements, tested against
\emph{literature-typical} rejoining curves, not against the paper's actual
figures or tables. Coverage of the paper's quantitative content is
approximately \textbf{0\%}. This is a NO-GO by data availability, not by
effort.

\section{Open Questions}
See \texttt{report/open\_questions.json} for machine-readable form.
\begin{description}
\item[Q1.] Does the DSB \emph{complexity distribution} (not just the total
number) explain the LET-scaling of misrejoining/mutation more than the
total DSB yield does? \emph{Basis:} Friedland 2012 encodes complexity
implicitly via the slow-channel saturation term but never publishes the
distribution shape parameters that PARTRAC generates. \emph{Next steps:}
(a) run Geant4-DNA at UICGPU for 5 ion species x 5 energies each, extract
per-DSB complexity histograms; (b) fit a two-component NHEJ model where
the slow rate is a function of DSB complexity bin (not aggregate LET); (c)
compare to KURBUC and RITRACKS on identical geometries.

\item[Q2.] How much does the NHEJ model output diverge across three
independent stochastic track-structure codes (PARTRAC vs Geant4-DNA vs
KURBUC) on the \emph{same} DSB spectrum input? \emph{Basis:} because
PARTRAC is closed, no independent group has been able to run a rigorous
cross-code test on Friedland's exact NHEJ formulation. The literature has
only paper-to-paper comparisons which conflate physics differences with
chemistry/biology differences. \emph{Next steps:} (a) implement the
Friedland 2012 NHEJ rules (fast/slow rejoining with complexity-dependent
capacity + labile-site delay) as a standalone Gillespie library, fed
externally with DSB end lists; (b) drive it with Geant4-DNA-generated DSB
lists, then with KURBUC-generated DSB lists, and diff the rejoining
kinetics; (c) publish the Gillespie library and its test-cases.

\item[Q3.] How sensitive are the fast/slow rate constants to \emph{chromatin
geometry} (fiber compaction, loop domain topology, hetero- vs
euchromatin), and does the ``two rate constants suffice'' claim survive
richer chromatin? \emph{Basis:} Friedland 2012 assumes a stylised chromatin
model within PARTRAC. Recent Hi-C-derived polymer models suggest DSB
proximity statistics depend strongly on chromatin state; if so, the
``labile-site delay'' Friedland attributes to biochemistry may be partly
geometric. \emph{Next steps:} (a) drive a Gillespie NHEJ model with DSB end
proximity graphs sampled from Hi-C-informed polymer simulations
(chromatin3D toolkit); (b) test whether fast/slow rate constants change by
$>$25\% between euchromatin- and heterochromatin-enriched regions.

\item[Q4.] What is the \emph{sub-DSB damage spectrum} (base damage, SSB
clusters, DSB with adjacent lesions) contribution to the apparent
``labile-site'' population, i.e.\ how much of the initial-phase detectable
DSB production is really delayed conversion of complex lesions rather than
enzyme kinetics? \emph{Basis:} Friedland 2012 phenomenologically encodes
labile-site delay but does not decompose it by lesion type. \emph{Next
steps:} (a) run Geant4-DNA with chemistry stage enabled to enumerate
sub-DSB lesion densities per LET; (b) fit a three-population model
(prompt-DSB / labile-converted DSB / clustered-lesion converted DSB) to
photon and ion kinetics; (c) test whether the labile amplitude scales with
predicted clustered lesion density.

\item[Q5.] Does the Friedland 2012 model extrapolate at all to the
\emph{FLASH ultra-high dose rate regime}, and if not, what minimal
extension is needed? \emph{Basis:} FLASH ($>$40 Gy/s) may alter DSB
production and repair via transient oxygen depletion or radical
recombination, phenomena absent from PARTRAC's conventional kinetics; the
paper predates FLASH and does not comment. \emph{Next steps:} (a)
literature-scan FLASH DSB rejoining measurements; (b) add a
dose-rate-dependent oxygen-depletion term to the Gillespie NHEJ library
and test whether the ``two-component + labile'' skeleton still fits FLASH
kinetics or whether a new fast$\to$very-fast component is needed; (c)
compare to Geant4-DNA-Chem simulations with radiolytic species.
\end{description}

\end{document}
